%!TEX TS-program = lualatex
%!TEX encoding = UTF-8 Unicode

\documentclass[11pt, addpoints]{exam}
\usepackage[left=1in,right=0.75in,top=1in]{geometry}     

\usepackage{fontspec}
\def\mainfont{Linux Libertine O}
\setmainfont[Ligatures=TeX, Contextuals={NoAlternate}, BoldFont={* Bold}, ItalicFont={* Italic}, BoldItalicFont={* BoldItalic}, Numbers={Proportional}]{\mainfont}
\setmonofont[Scale=MatchLowercase]{Inconsolata} 
\setsansfont[Scale=MatchLowercase]{Linux Biolinum O} 
\usepackage{microtype}

\usepackage{graphicx}
	\graphicspath{%
	{/Users/goby/Pictures/teach/151/exams/}%
	{/img/}}%

% Used to get the semestesr and last two digits of the year for the header below.
\def\short#1{\csname @gobbletwo\expandafter\endcsname\number#1}
\def\Semester{\ifthenelse{\number\month>8}
	{F}%
	{%
		\ifthenelse{\number\month<4}
		{S}%
		{N}%
	}%
}

\pagestyle{headandfoot}
\firstpageheader{BI 151, Quizzam 2, \Semester{}\short{\year}}{}{%
	\ifprintanswers\textbf{KEY}\else Name: \enspace \makebox[2.5in]{\hrulefill}\fi}
\runningheader{}{}{\small{pg. \thepage}}

\footer{}{}{}
\extrafootheight{-0.5in}

\pointsinmargin
\marginpointname{ pts}

\usepackage{multicol}
\usepackage{booktabs}

%% Remove comment %% to print answer key.
%\printanswers
\unframedsolutions
\renewcommand{\solutiontitle}{}
\SolutionEmphasis{\bfseries}

%% Create a Matching question format
\newcommand*\Matching[1]{
\ifprintanswers
	\textbf{#1}
\else
	\rule{2.1in}{0.5pt}
\fi
}
\newlength\matchlena
\newlength\matchlenb
\settowidth\matchlena{\rule{2.1in}{0pt}}
\newcommand\MatchQuestion[2]{%
	\setlength\matchlenb{\linewidth}
	\addtolength\matchlenb{-\matchlena}
	\parbox[t]{\matchlena}{\Matching{#1}}\enspace\parbox[t]{\matchlenb}{#2}}



\begin{document}

\fullwidth{Read each question completely before answering. Write clear and concise answers. \vspace{2\baselineskip}}

\begin{questions}

\question[4]
Here is the full classification of the Scalloped Hammerhead Shark. 

Eukaryota; Metazoa; Chordata; Chondrichthys; Carcharhiniformes; Sphyrnidae; Sphyrna; lewini \vspace{\baselineskip}

\begin{parts}

	\part Which is the order? 
	\ifprintanswers \textbf{Charcharhiniformes}\\\vspace{\baselineskip}\else\vspace{2\baselineskip}\fi
	
	\part Write the scientific name in proper format.
	\ifprintanswers \textbf{\textit{Sphyrna lewini}}\\\vspace{\baselineskip}\else\vspace{2\baselineskip}\fi
	
\end{parts}

%\vspace*{\stretch{1}}

\question[2]
The type of scientific reasoning that uses a series of specific observations to develop general principles is called:
	\begin{solution}
	Inductive reasoning.
	\end{solution}	

\vspace*{\stretch{0.25}}

\question[4]
Clearly explain the difference between a belief and a scientific theory.

	\begin{solution}
	A belief is something accepted to be true without evidence. A theory is a framework of scientific hypotheses that has been well studied and explains facts about the natural world.
	\end{solution}

\vspace*{\stretch{1}}

\question[4]
Write a clear definition of homology.

	\begin{solution}
	A homology is a structural similarity between two or more organisms that is due to common ancestry.
	\end{solution}
	
\vspace*{\stretch{1}}

\newpage

\fullwidth{Consider the following hypothesis about the relationships among various organisms.
\begin{center}
	\includegraphics[width=0.95\textwidth]{exam1_figure1}
\end{center}

Answer questions \ref{first_phylo}–\ref{last_phylo} based on the hypothesis above. \emph{Tell how you know for each answer!}}

\question[3]\label{first_phylo}
Which two organisms shared a common ancestor \emph{most recently}? 

	\begin{solution}
	Nilgiri Marten and Pine Marten. They arose about 5 MYA. The next closest splits are about 10 MYA.
	\end{solution}

\vspace*{\stretch{1}}
	
\question[3]
Which are most closely related, Nilgiri marten and sable, or carpenter bee and bumble bee?

	\begin{solution}
		Bumble bee and carpenter bee share the same common ancestor. 
		Nilgiri marten and Sable trace back one additional common ancestor.
	\end{solution}
\vspace*{\stretch{1}}

\question[3]
Which organism(s), if any, is (are) most closely related to a hornet?
	\begin{solution}
		Mud dauber and Paper wasp are most closely related. They share the most recent common ancestor with hornet.
	\end{solution}
\vspace*{\stretch{1}}

\newpage

\question[3]
Approximately when did the common ancestor to big-eared kangaroo rat and desert kangaroo rat first appear?

	\begin{solution}
		About 30 million years ago. Their common ancestor ends at that time and they begin at that time.
	\end{solution}
\vspace*{\stretch{1}}

\question[3]
Which organism(s), if any, is most closely related to the bobcat?
	\begin{solution}
		None. This hypothesis predicts that the bobcat does not have any common ancestors.
	\end{solution}


\vspace*{\stretch{1}}

\question[3]
Suppose you find a homology between bumble bee and mud dauber. Is the
hypothesis supported, falsified, or is the evidence inconclusive?
	\begin{solution}
		The hypothesis predicts the bumble bee and mud dauber have a common ancestor and therefore should have a homology. If a homology is found, the hypothesis is supported.
	\end{solution}

\vspace*{\stretch{1}}

\question[3]
Suppose you find an analogy between desert kangaroo rat and bobcat. Is the
hypothesis supported, falsified, or is the evidence inconclusive?
	\begin{solution}
		Analogies do not provide evidence that can be used to test hypothesis. The evidence is inconclusive. 
	\end{solution}

\vspace*{\stretch{1}}

\question[3]\label{last_phylo}
Suppose you find a homology between big-eared kangaroo rat and pine marten. Is
the hypothesis supported, falsified, or is the evidence inconclusive?
	\begin{solution}
		The hypothesis predicts that big-eared kangaroo rat and pine marten do not have a common ancestor. Therefore, they should not have a homology. If a homology is found, then the hypothesis is falsified.
	\end{solution}
\vspace*{\stretch{1}}

\newpage

\question[7]
List below the four missing steps of the Scientific Method. Then,
for each step, write a statement that illustrates how scientists applied
the scientific method to the experiments described below. I have
provided step four as an example.

\textbf{Houseflies are infamous for walking all over food around the house
and at picnics. They stick out their proboscis (the fly version of a tongue)
when they walk on food but not when the walk on other things, like the table.
Scientists thought that flies might be able to taste the food with their feet,
and stick out their proboscis whenever they identified something tasty. To
test this idea, scientists filled one saucer with sugar water (sugar is a food
source for the fly) and one saucer with plain water. They then glued live 
flies to the end of a slender stick. The flies were gently lowered until their 
feet just touched the liquid. If the liquid was plain water, the flies did not stick
out their proboscis. If the liquid was sugar water, the flies always stuck out 
their proboscis.}

%The Five Steps of the Scientific Method:

\begin{parts}
	\part	\vspace{\stretch{0.25}} 
	\ifprintanswers
		\textbf{Observation:} Flies stick out their tongue when walking on food.
	\fi
	
	\part \vspace{\stretch{1}}
	\ifprintanswers
		\textbf{Hypothesis:} Flies taste with their feet.
	\fi
	
	\part \vspace{\stretch{1}}
	\ifprintanswers
		\textbf{Prediction:} If fly feet are touched to a food source, they will stick out their proboscis.
	\fi

	\part \vspace{\stretch{1}} \textbf{Experiment and results}: Scientists noted that flies stuck out
	their proboscis when their feet touched sugar water but not when their feet touched 
	plain water.	

	\part \vspace*{\stretch{1}}
	\ifprintanswers
		\textbf{Conclusion:} The hypothesis was supported.
	\fi
\end{parts}

	\vspace*{\stretch{1}}

\newpage

\question[5]
Below are several \emph{correct} ways to draw phylogenetic trees. Which tree shows a different pattern of \emph{relationships} compared to the others. Identify \emph{all} differences between the tree you identified and the other three trees.  Tell how you know why the relationships are different.

	\includegraphics[width=0.95\textwidth]{exam1_figure2}

	\begin{solution}
		Tree C.  Tree C shows D and G sharing a most recent common ancestor and E and F sharing a most recent common ancestor. The other three trees show D and E sharing the most recent common ancestor. G shares a common ancestor with F in the other trees.\vspace{\baselineskip}
		
		1 pt for correct tree, 2 pts for explaining each of the two differences.
	\end{solution}
	
	
\end{questions}


\end{document}